% Content of the blueprint for the 3-majority opinion-dynamics formalization.
% Mirrors latex/three_majority.tex in spirit, but every statement here also
% carries \lean{...} / \leanok / \uses{...} tags matched to the *actual*
% Lean declarations in ThreeMajority/*.lean (namespace `ThreeMajority`), so
% that leanblueprint can build the dependency graph and check coverage.
% Where the Lean proof takes a different route than the paper draft in
% latex/three_majority.tex, this file follows the Lean proof.

\begin{abstract}
\noindent
We give a complete, elementary proof that the \emph{3-majority}
opinion-dynamics process on $n$ fully-mixing agents, each holding one of two
opinions, reaches consensus within $O(\log n)$ rounds with probability
$1-O(1/n)$, provided the initial imbalance between the two opinions exceeds
a fixed constant fraction of $n$ (at least $60\%$ vs.\ $40\%$). The
probability space is finite and uniform throughout; the only analytic
ingredient beyond finite sums is a single self-contained exponential-moment
(Chernoff-type) concentration lemma for sums of independent
$\{0,1\}$-valued random variables, proved from the elementary bound
$1+x \le e^x$ --- no measure theory, no PMF/kernels, no appeal to Mathlib's
probability library. This is a genuinely different regime from the
companion \texttt{rumor\_spread} (push protocol) development: the
opinion-$1$ count is \emph{not} monotone in the round index (an unlucky
round can shrink it), so a counting argument over ``good rounds'' does not
suffice --- honest concentration around the process's mean trajectory is
needed in \emph{both} the growth and saturation phases. This argument has
been \emph{fully formalized} (no \texttt{sorry}) in Lean~4 + Mathlib; the
main theorem is \texttt{ThreeMajority.majority3\_consensus\_whp}.

Section~\ref{sec:bcekmn-any-start} formalizes, for any number of colours, the comparison of
3-Majority with the Voter model of Berenbrink et al.\ (PODC 2017) and their bound
$O(n^{3/4}\log^{7/8}n)$ on the consensus time from any configuration. Its second phase is
Theorem~3.1 of Becchetti et al.\ (SODA 2016), from $k\le n^{1/3-\varepsilon}$ colours, which
Section~\ref{sec:bcnpt16-many-opinions} formalizes; together they give the bound without
hypothesis.
\end{abstract}

\section{The model}\label{sec:model}

Fix $n \ge 1$ and let $V = \{0,1,\dots,n-1\}$ be the set of agents, each
holding one of two opinions, $1$ (``red'') or $0$ (``blue''). Write
$I \subseteq V$ for the set of agents currently holding opinion $1$.

\begin{definition}[Round randomness]\label{def:round}
  \lean{ThreeMajority.Tgt3}
  \leanok
  A \emph{round configuration} is a function $r \colon V \to V \times V
  \times V$ assigning to every agent $v$ an ordered triple of \emph{samples},
  each drawn independently and uniformly from $V$, \emph{with} replacement
  (an agent may sample itself, or the same agent twice --- harmless, since a
  majority of three Boolean values is always well defined). Let
  $R_n = V^{3V}$ denote the (finite) set of round configurations, with the
  uniform probability measure.
\end{definition}

\begin{definition}[Majority step]\label{def:step}
  \lean{ThreeMajority.step, ThreeMajority.sampleCount, ThreeMajority.sampleCountOf, ThreeMajority.mem_step}
  \leanok
  \uses{def:round}
  For an opinion-$1$ set $I \subseteq V$ and $r \in R_n$,
  \[
    \mathrm{step}(I,r) \;=\; \{\, v \in V : \text{at least two of } r(v)\text{'s three
    coordinates lie in } I \,\}.
  \]
  Each agent independently samples $3$ agents and adopts the majority of
  their current opinions.
\end{definition}

\begin{definition}[Process]\label{def:process}
  \lean{ThreeMajority.run}
  \leanok
  \uses{def:step}
  Fix a horizon $T \in \N$ and an initial set $I_0 \subseteq V$. The sample
  space is $\Omega = R_n^{\,T}$ with the uniform (i.i.d.) measure. For
  $r = (r_0,\dots,r_{T-1}) \in \Omega$ define $I_{t+1} =
  \mathrm{step}(I_t, r_t)$. We write $X_t = |I_t|$ and $U_t = n - X_t$. In
  the formalization, a trajectory is a \texttt{List} of round configurations
  and $\mathrm{run}\ I\ l$ folds \texttt{step} over the list $l$.
\end{definition}

\begin{lemma}[Monotone coupling]\label{lem:mono}
  \lean{ThreeMajority.step_mono, ThreeMajority.run_mono, ThreeMajority.sampleCount_mono}
  \leanok
  \uses{def:step, def:process}
  For a \emph{fixed} round $r$, the map $I \mapsto \mathrm{step}(I,r)$ is
  monotone: $I \subseteq I'$ implies $\mathrm{step}(I,r) \subseteq
  \mathrm{step}(I',r)$; iterating over a trajectory, $\mathrm{run}$ is
  monotone in the same sense. This is purely combinatorial (no probability
  involved).
\end{lemma}

\begin{remark}
  Unlike the push protocol, where the analogous monotonicity
  (Lemma~\ref{lem:mono} there) is load-bearing throughout, here the
  per-round probabilistic estimates below (Lemmas~\ref{lem:growthround},
  \ref{lem:satclosed}, etc.) are instead stated and proved directly for
  \emph{any} set $I$ satisfying a hypothesis on $|I|$, using monotonicity of
  the real polynomial $p$ of Section~\ref{sec:map} on $[0,1]$ rather than a
  reference-embedding argument via Lemma~\ref{lem:mono}. Lemma~\ref{lem:mono}
  is recorded because it is the exact analogue of the push protocol's
  structural fact, but it is not otherwise used by the chain of lemmas
  leading to the main theorem.
\end{remark}

\section{Main theorem}\label{sec:main-statement}

Throughout, $\log$ denotes the natural logarithm.

\begin{theorem}[3-majority reaches consensus in $O(\log n)$ rounds w.h.p.]
  \label{thm:main}
  \lean{ThreeMajority.majority3_consensus_whp}
  \leanok
  \uses{lem:mainfailclean}
  Let $n$ be such that $\log n \ge 30$ (in particular $n \ge 2$), let
  $I_0 \subseteq V$ with $|I_0| \ge \tfrac35 n$, and set
  \[
    T \;=\; 10 + (T_{2a}(n) + 2), \qquad T_{2a}(n) = \ceil{6\log n}.
  \]
  Then $\Pr[\, I_T = V \,] \;\ge\; 1 - \dfrac{500}{n}$.
\end{theorem}

\begin{remark}
  As in the push protocol's write-up, the threshold $\log n \ge 30$ and
  the constant $500$ are deliberately crude: the proof is organized to keep
  every inequality an easy direct computation, not to be sharp. Reaching
  this modest, two-digit threshold on $\log n$ (rather than an
  astronomically large one) relies on two ingredients in
  Section~\ref{sec:ineq}: a degree-$11$ instance of the generic
  exponential-vs-power bound (Lemma~\ref{lem:exppow}), and a log-log bound
  tangent at a reference point near the threshold itself rather than at the
  fixed point $e$ (Lemma~\ref{lem:loglog}). It is known
  (Doerr--Goldberg--Minder--Sauerwald--Scheideler 2011;
  Becchetti--Clementi--Natale--Pasquale--Silvestri, SODA 2015) that an
  initial imbalance of only $\Theta(\sqrt{n\log n})$ already suffices for
  $O(\log n)$-round consensus w.h.p.; reaching that sharp threshold
  requires tracking the process on the scale of its own standard deviation
  throughout, a substantially more delicate multi-scale argument. Here a
  constant-fraction initial imbalance keeps every step to a single,
  reusable Chernoff bound (Section~\ref{sec:chernoff}) applied at a handful
  of explicit parameters.
\end{remark}

The proof has two phases, as in the push protocol, but for a structurally
different reason: \textbf{growth} (Section~\ref{sec:growth}), where the
opinion-$1$ fraction self-amplifies from $60\%$ to $75\%$; and
\textbf{saturation} (Section~\ref{sec:saturation}), where the last
dissenting agents are eliminated in three stages. Because the process is
not monotone in the round index, \emph{both} phases need the Chernoff bound
of Section~\ref{sec:chernoff} --- unlike the push protocol, where only
saturation needed a (much weaker) expectation argument.

\section{Finite uniform probability and independence}\label{sec:prelim}

As in the push protocol, all randomness is uniform over finite types, so
instead of measure theory or \texttt{PMF} we work with plain finite sums.

\begin{definition}[Average]\label{def:avg}
  \lean{ThreeMajority.avg}
  \leanok
  For a finite type $\alpha$ and $f \colon \alpha \to \R$,
  $\mathrm{avg}(f) := \bigl(\sum_a f(a)\bigr)/|\alpha|$ is the expectation of
  $f$ under the uniform distribution on $\alpha$.
\end{definition}

\begin{definition}[Trajectory expectation]\label{def:expList}
  \lean{ThreeMajority.expList}
  \leanok
  \uses{def:avg}
  Exactly as in the push protocol: $\mathrm{expList}\ \alpha\ T\ F$ is the
  expectation of a functional $F$ of a trajectory of $T$ i.i.d.\ uniform
  draws from $\alpha$, defined by recursion on $T$ (the head draw is
  averaged out first), so conditioning on the first round is a definitional
  unfolding.
\end{definition}

\begin{lemma}[Basic properties]\label{lem:avg-basic}
  \lean{ThreeMajority.avg_nonneg, ThreeMajority.avg_le_avg, ThreeMajority.avg_add,
    ThreeMajority.avg_sub, ThreeMajority.avg_const_mul, ThreeMajority.avg_sum,
    ThreeMajority.avg_indicator, ThreeMajority.avg_const,
    ThreeMajority.expList_nonneg, ThreeMajority.expList_le_expList,
    ThreeMajority.expList_add, ThreeMajority.expList_const_mul,
    ThreeMajority.expList_const, ThreeMajority.expList_append}
  \leanok
  \uses{def:avg, def:expList}
  $\mathrm{avg}$ and $\mathrm{expList}$ are monotone, additive, homogeneous,
  and constants pass through unchanged; $\mathrm{expList}$ over a
  concatenated horizon splits as a nested $\mathrm{expList}$
  (\texttt{expList\_append}), used to glue phases in
  Section~\ref{sec:proof-main}.
\end{lemma}

\begin{lemma}[Independence over a product space]\label{lem:indep}
  \lean{ThreeMajority.avg_mul_prod, ThreeMajority.avg_fst_mul, ThreeMajority.avg_snd_mul,
    ThreeMajority.avg_equiv, ThreeMajority.avg_prod_pi, ThreeMajority.avg_eval}
  \leanok
  \uses{def:avg}
  \textbf{New relative to the push protocol}, since that development only
  ever averaged over a single uniform draw at a time. For finite types
  $\beta,\delta$ and $g \colon \beta\to\R$, $h\colon\delta\to\R$:
  $\mathrm{avg}(p \mapsto g(p_1)h(p_2)) = \mathrm{avg}(g)\,\mathrm{avg}(h)$
  over the product $\beta\times\delta$ (\texttt{avg\_mul\_prod}), with
  marginal consequences \texttt{avg\_fst\_mul}/\texttt{avg\_snd\_mul}.
  Iterating over a $\mathrm{Fin}\ n$-indexed product
  (\texttt{avg\_prod\_pi}): for $f \colon \mathrm{Fin}\ n \to \gamma \to \R$,
  \[
    \mathrm{avg}\Bigl(x \mapsto \textstyle\prod_i f_i(x_i)\Bigr)
      \;=\; \prod_i \mathrm{avg}(f_i)
  \]
  over $x \colon \mathrm{Fin}\ n \to \gamma$ --- the discrete,
  measure-theory-free form of ``coordinate projections on a product space
  are independent,'' with a one-coordinate marginal corollary
  (\texttt{avg\_eval}). This is exactly the independence fact that lets the
  $n$ agents' per-round updates be combined into a single concentration
  bound (Section~\ref{sec:chernoff}).
\end{lemma}

\section{The majority map}\label{sec:map}

Write $x = X_t/n \in [0,1]$ for the opinion-$1$ fraction, and
$\mathrm{ind}(I)$ for the real-valued membership indicator of $I$, so
$\mathrm{avg}(\mathrm{ind}(I)) = |I|/n$.

\begin{lemma}[Majority-of-three as a polynomial]\label{lem:majpoly}
  \lean{ThreeMajority.ind, ThreeMajority.avg_ind, ThreeMajority.avg_ind_eq,
    ThreeMajority.majorityIndicator_eq}
  \leanok
  \uses{def:avg, def:step}
  For $a,b,c \in \{0,1\}$: $\mathbf 1[a+b+c\ge2] = ab+bc+ac-2abc$. Applied
  pointwise to the three samples of a fixed agent, this expresses the
  majority indicator as a polynomial in the three membership indicators.
\end{lemma}

\begin{definition}[Per-agent contribution]\label{def:Ymaj}
  \lean{ThreeMajority.Y_maj, ThreeMajority.Y_maj_zero_one, ThreeMajority.card_step_eq_sum}
  \leanok
  \uses{lem:majpoly}
  $Y_{\mathrm{maj}}(I,v,s) \in \{0,1\}$ is $1$ iff a majority of the sample
  triple $s$ lies in $I$ (independent of $v$, kept as a function of $v$ to
  match the shape the Chernoff bound of Section~\ref{sec:chernoff} expects).
  Then $|\mathrm{step}(I,r)| = \sum_v Y_{\mathrm{maj}}(I,v,r(v))$
  (\texttt{card\_step\_eq\_sum}): the per-agent, $\{0,1\}$-valued
  decomposition that Section~\ref{sec:chernoff}'s tail bounds need.
\end{definition}

\begin{lemma}[Exact one-round drift]\label{lem:drift}
  \lean{ThreeMajority.avg_ind_mul_fst_snd1, ThreeMajority.avg_ind_mul_snd1_snd2,
    ThreeMajority.avg_ind_mul_fst_snd2, ThreeMajority.avg_ind_mul_triple,
    ThreeMajority.avg_Y_maj_eq, ThreeMajority.sum_avg_Y_maj, ThreeMajority.avg_card_step}
  \leanok
  \uses{def:Ymaj, lem:indep}
  Conditionally on a fixed $I$ with $|I|=m$, the next count
  $X_{t+1} = |\mathrm{step}(I,r)|$ satisfies
  \[
    \E_r[X_{t+1}] \;=\; n\,p(x), \qquad p(x) := 3x^2-2x^3, \qquad x = m/n,
  \]
  proved exactly (no approximation) via Lemma~\ref{lem:majpoly},
  Definition~\ref{def:Ymaj} and the independence toolkit of
  Lemma~\ref{lem:indep} (each pairwise/triple product of membership
  indicators over distinct coordinates of the product space
  $\mathrm{Tgt3}\ n$ factors by \texttt{avg\_mul\_prod}).
\end{lemma}

\begin{definition}[Dissent contribution]\label{def:Ydis}
  \lean{ThreeMajority.Y_dis, ThreeMajority.Y_dis_zero_one, ThreeMajority.card_dissent_eq_sum,
    ThreeMajority.sum_avg_Y_dis}
  \leanok
  \uses{def:Ymaj, lem:drift}
  $Y_{\mathrm{dis}} := 1 - Y_{\mathrm{maj}}$, the complementary per-agent
  dissent contribution, so $n - |\mathrm{step}(I,r)| = \sum_v
  Y_{\mathrm{dis}}(I,v,r(v))$ with mean $n - n\,p(x)$
  (\texttt{sum\_avg\_Y\_dis}); by the algebraic identity
  $1-p(x) = (1-x)^2(1+2x)$ this is the quantity controlled in
  Section~\ref{sec:saturation}.
\end{definition}

\begin{remark}
  The algebraic identities for $p$ used throughout below --- $p(x)-x =
  x(1-x)(2x-1)$, $1-p(x) = (1-x)^2(1+2x)$, and the growth/saturation
  regional bounds $p(\tfrac12+\beta)-\tfrac12 \ge \tfrac54\beta$ for
  $0\le\beta\le\tfrac14$ and $u(3-2u)\le\tfrac58$ for $0\le u\le\tfrac14$ ---
  are not isolated as separate Lean lemmas: each use site
  (\texttt{growth\_round}, \texttt{saturation\_round\_closed},
  \texttt{saturation\_mean\_quad}) re-derives the specific polynomial
  inequality it needs inline via \texttt{ring}/\texttt{nlinarith}, in the
  same spirit as the push protocol's reverse-Markov bound, which is also
  not isolated as its own named lemma.
\end{remark}

\section{Elementary inequalities}\label{sec:ineq}

\begin{lemma}[Quadratic-tight logarithm bound]\label{lem:logquad}
  \lean{ThreeMajority.log_ge_quadratic}
  \leanok
  For $\tfrac12 \le x \le 1$: $(x-1)-(x-1)^2 \le \log x$. (Mathlib's crude
  linear bound $\log x \ge x-1$ is false for $x<1$ and, worse, plugging the
  correct one-sided bound $1-\tfrac1x\le\log x$ into the Chernoff exponent
  $k-\mu-k\log(k/\mu)$ gives exactly $0$ when $k/\mu\to1$ --- useless. This
  quadratic correction, derived from Mathlib's degree-$3$ Taylor lower bound
  for $\exp$ by elementary algebra, is what makes the growth-phase Chernoff
  bound (Lemma~\ref{lem:growthround}) bite.)
\end{lemma}

\begin{lemma}[Exponential beats any fixed power]\label{lem:exppow}
  \lean{Real.pow_div_factorial_le_exp, ThreeMajority.exp_ge_cube}
  \leanok
  For $x\ge0$ and any $k\in\N$: $x^k/k! \le \exp(x)$ --- the single
  degree-$k$ term of the Taylor truncation, extracted by dropping every
  other (nonnegative, since $x\ge0$) term of the sum. The cubic case
  $k=3$ (\texttt{exp\_ge\_cube}) is the special case used by
  Lemma~\ref{lem:logquad}. Choosing $k$ large enough relative to a given
  lower bound on $\log n$ lets $n$ be shown to dominate any fixed
  polynomial in $\log n$ using only a \emph{modest} such lower bound
  (Stage 2b, Lemma~\ref{lem:stage2b}, and the final numeric bound
  Lemma~\ref{lem:mainfailclean}) --- a single low-degree truncation, such
  as the cubic case alone, would instead force $\log n$ to be
  astronomically large.
\end{lemma}

\begin{lemma}[Tangent-line log bound]\label{lem:logtangent}
  \lean{ThreeMajority.log_le_tangent_div, ThreeMajority.exp_three_ge_twenty}
  \leanok
  For $L>0$ and any reference point $c>0$: $\log L \le \log c - 1 + L/c$,
  from the crude bound $\log t \le t - 1$ applied to $t = L/c$; tight at
  $L=c$ (equality there), unlike a bound tangent at a single globally-fixed
  point, so choosing $c$ near the $L$ actually in play gives a far sharper
  bound. Also records $\exp(3)\ge20$ (cubing the standard $10$-digit lower
  bound on $e$), used to instantiate $c=\exp(3)$ below.
\end{lemma}

\begin{lemma}[Log-log bound]\label{lem:loglog}
  \lean{ThreeMajority.log_log_le_of_pos}
  \leanok
  \uses{lem:logtangent}
  For $\log x>0$: $\log\log x \le 2 + (\log x)/20$, from
  Lemma~\ref{lem:logtangent} tangent at the reference point
  $c=\exp(3)\approx20.09$ (so $\log c = 3$ exactly), combined with
  $\exp(3)\ge20$ to replace $1/\exp(3)$ by the cruder but rational $1/20$.
  Far tighter, once $\log x$ is more than a few units above $e$, than a
  bound tangent at the fixed point $c=e$ (the crude bound $\log t\le t-1$
  applied directly at $t=(\log x)/e$) --- exactly the regime
  Lemma~\ref{lem:stage2b} needs. Used only inside Lemma~\ref{lem:stage2b}'s
  numeric estimate.
\end{lemma}

\section{An elementary Chernoff bound}\label{sec:chernoff}

The only concentration tool used in the whole development, and the only
place beyond elementary combinatorics where real exponentials appear. These
bounds are proved once in the shared \texttt{dynamics} library
(\texttt{Dynamics/Chernoff.lean}, namespace \texttt{Dynamics}), which the
plurality and median developments use as well.

\begin{lemma}[MGF bound]\label{lem:mgf}
  \lean{Dynamics.avg_chernoff_mgf}
  \leanok
  \uses{lem:indep}
  Let $Y_1,\dots,Y_n \colon \gamma \to \{0,1\}$ be the per-agent
  contributions on the product space $\mathrm{Fin}\ n \to \gamma$ (one
  independent draw of $\gamma$ per agent), with $\mu=\sum_i\mathrm{avg}(Y_i)$
  and $X(x) := \sum_i Y_i(x_i)$. Then for every real $t$ (no sign
  restriction):
  \[
    \E[\exp(tX)] \;\le\; \exp\bigl(\mu(e^t-1)\bigr).
  \]
  Proved via the independence toolkit (Lemma~\ref{lem:indep}) applied to
  $x \mapsto \prod_i \exp(tY_i(x_i))$, and the elementary bound
  $1+y\le e^y$ (Mathlib's \texttt{Real.add\_one\_le\_exp}) at each factor.
\end{lemma}

\begin{lemma}[Chernoff tail bounds]\label{lem:chernoff}
  \lean{Dynamics.avg_chernoff_upper_of_mgf, Dynamics.avg_chernoff_lower_of_mgf}
  \leanok
  \uses{lem:mgf}
  In the setting of Lemma~\ref{lem:mgf}: for $t\ge0$ and any $k$,
  $\Pr[X\ge k] \le \exp(\mu(e^t-1)-tk)$ (\texttt{avg\_chernoff\_upper\_of\_mgf}, Markov
  applied to $\exp(tX)$); symmetrically for $t\le0$,
  $\Pr[X\le k] \le \exp(\mu(e^t-1)-tk)$ (\texttt{avg\_chernoff\_lower\_of\_mgf}).
\end{lemma}

\begin{lemma}[Closed-form tail bounds at the optimal $t$]\label{lem:chernofflog}
  \lean{Dynamics.avg_chernoff_upper_log, Dynamics.avg_chernoff_lower_log}
  \leanok
  \uses{lem:chernoff}
  Plugging $t=\log(k/\mu)$ into Lemma~\ref{lem:chernoff}: for $k\ge\mu>0$,
  \[
    \Pr[X\ge k] \;\le\; \exp\bigl(k-\mu-k\log(k/\mu)\bigr)
    \qquad(\mu=\E X),
  \]
  and symmetrically for $0<k\le\mu$,
  $\Pr[X\le k] \le \exp(k-\mu-k\log(k/\mu))$.
  Both closed forms are \emph{mean-scaled}: the exponent depends on the true
  $\mu$ directly, not on a crude worst-case range/variance proxy, which
  matters late in the saturation phase once few at-risk agents remain
  (Lemma~\ref{lem:stage2b}). The Lean statements are the generalized
  variants, which take only an \emph{upper} bound $\mu_{\mathrm{ub}}\ge\mu$
  (\texttt{avg\_chernoff\_upper\_log}, used by the saturation phase) or a
  \emph{lower} bound $\mu_{\mathrm{lb}}\le\mu$ (\texttt{avg\_chernoff\_lower\_log}, used by
  the growth phase) on the true mean (the exact forms are the case of equality) --- monotonicity of the exponent in $\mu$ at the optimal $t$
  makes both directions sound, and this is the form each phase actually
  needs, since the exact mean depends on $x=|I|/n$ while the phases only
  assume an inequality on $|I|$.
\end{lemma}

\section{Growth phase}\label{sec:growth}

Set $\beta_j = \tfrac1{10}(1.1)^j$ for $j=0,\dots,10$; $\beta_0=0.1$ (the
hypothesis $X_0\ge0.6n$) and $\beta_{10}>0.25$.

\begin{definition}[Bias schedule]\label{def:growthbias}
  \lean{ThreeMajority.growthBias, ThreeMajority.growthBias_zero, ThreeMajority.growthBias_nonneg,
    ThreeMajority.growthBias_mono, ThreeMajority.growthBias_le_of_le,
    ThreeMajority.growthBias_le_quarter, ThreeMajority.growthBias_ten}
  \leanok
  $\mathrm{growthBias}(j) := \tfrac1{10}(\tfrac{11}{10})^j$, the deterministic
  bias target after $j$ growth rounds; monotone, $\le\tfrac14$ for $j\le9$,
  and $>\tfrac14$ at $j=10$.
\end{definition}

\begin{lemma}[Single growth round]\label{lem:growthround}
  \lean{ThreeMajority.growth_round}
  \leanok
  \uses{lem:drift, lem:chernofflog, def:growthbias}
  Starting from opinion-$1$ fraction $\ge\tfrac12+\beta$
  ($\tfrac1{10}\le\beta\le\tfrac14$), the probability of \emph{failing} to
  reach the next bias target (guaranteed growth factor $\tfrac{11}{10}$,
  comfortably below the true drift factor $\tfrac54$ from
  Lemma~\ref{lem:drift}) is at most $\exp(-n/10^7)$. Proved by computing a
  lower bound $\mu_{\mathrm{lb}} = n(\tfrac12+\tfrac54\beta)$ on the true
  mean via monotonicity of $p$ on $[\tfrac12+\beta,1]$ (an inline
  \texttt{nlinarith} computation, not a separate lemma), then applying the
  lower-tail bound \texttt{avg\_tail\_le\_log\_ge} at
  $k=n(\tfrac12+\tfrac{11}{10}\beta)$, and bounding the resulting exponent
  using Lemma~\ref{lem:logquad}.
\end{lemma}

\begin{lemma}[Union bound over growth rounds]\label{lem:growthfail}
  \lean{ThreeMajority.growth_fail_le}
  \leanok
  \uses{lem:growthround}
  Starting from bias $\mathrm{growthBias}(i)$, after $j$ further rounds
  (with $i+j\le10$) the probability of failing to reach
  $\mathrm{growthBias}(i+j)$ is at most $j\exp(-n/10^7)$, by induction on
  $j$ via the $\mathrm{expList}$-conditioning idiom (mirroring the push
  protocol's \texttt{Main.lean}).
\end{lemma}

\begin{lemma}[Phase 1]\label{lem:phase1}
  \lean{ThreeMajority.growth_phase1}
  \leanok
  \uses{lem:growthfail}
  Starting from opinion-$1$ fraction $\ge\tfrac35$, after $10$ rounds the
  opinion-$1$ fraction exceeds $\tfrac34$, except with probability
  $\le10\exp(-n/10^7)$.
\end{lemma}

\section{Saturation phase}\label{sec:saturation}

Write $U_t = n-X_t$. By Lemma~\ref{lem:phase1}, $U_{10}\le n/4$ except with
probability $\le10\exp(-n/10^7)$. Getting $U_t$ all the way to
\emph{exactly} $0$ needs three stages, since a fixed-ratio Chernoff bound
needs the mean itself to be $\Omega(\log n)$ for a polynomially-small
per-round failure probability.

\begin{lemma}[General per-round upper-tail bound]\label{lem:satclosed}
  \lean{ThreeMajority.saturation_round_closed, ThreeMajority.saturation_round_generic}
  \leanok
  \uses{def:Ydis, lem:chernofflog}
  Given a reference dissent bound $U_t\le M\le n/4$ and any target
  $k\ge\tfrac58M$ (an upper bound $\mu_{\mathrm{ub}}=\tfrac58M$ on the true
  one-round mean, via $1-p(x)=(1-x)^2(1+2x)\le\tfrac58(1-x)$ for
  $x\ge\tfrac34$), the probability the next dissent count exceeds $k$ is at
  most $\exp(k-\tfrac58M-k\log(k/(\tfrac58M)))$
  (\texttt{saturation\_round\_closed}); the same conclusion holds for any
  directly-supplied upper bound $\mu_{\mathrm{ub}}$ on the true mean
  (\texttt{saturation\_round\_generic}, used by Stage 2b where the
  quadratic estimate below is the tight one).
\end{lemma}

\begin{lemma}[Quadratic mean bound]\label{lem:satquad}
  \lean{ThreeMajority.saturation_mean_quad}
  \leanok
  \uses{def:Ydis}
  For any dissent bound $U_t\le M\le n$, the true one-round mean is at most
  $3M^2/n$ (the crude bound $1-p(x)\le3(1-x)^2$, dropping the $\tfrac58$
  refinement). Used both by Stage 2b and by Stage 2c (at $M=10$).
\end{lemma}

\begin{lemma}[Numeric per-round bound]\label{lem:satnumeric}
  \lean{ThreeMajority.saturation_round_numeric}
  \leanok
  \uses{lem:satclosed, lem:logquad}
  Lemma~\ref{lem:satclosed}'s exponent, simplified via Mathlib's
  quadratic-tight Pad\'e bound for $\log$
  (\texttt{Real.le\_log\_one\_add\_of\_nonneg}) to the closed form
  $-\lambda^2/(k+\tfrac58M)$, $\lambda=k-\tfrac58M$.
\end{lemma}

\begin{lemma}[Stage 2a per-round contraction]\label{lem:satcontract}
  \lean{ThreeMajority.saturation_round_contract}
  \leanok
  \uses{lem:satnumeric}
  Guaranteed contraction factor $0.7$ (vs.\ the true $\tfrac58$): failure
  probability $\le\exp(-M/250)$.
\end{lemma}

\begin{definition}[Deterministic dissent target]\label{def:satafter}
  \lean{ThreeMajority.satAfter, ThreeMajority.satAfter_ge_floor, ThreeMajority.satAfter_zero,
    ThreeMajority.satAfter_le_of_le, ThreeMajority.satAfter_contract_le}
  \leanok
  $\mathrm{satAfter}(n,M,i) := \max(500\log n,\ (0.7)^i M)$, the target
  dissent bound after $i$ Stage 2a rounds, clamped below at the floor
  $500\log n$; written in closed form (not recursively) so shifts by one
  round agree definitionally on the unclamped branch.
\end{definition}

\begin{lemma}[Stage 2a, union bound]\label{lem:satfail}
  \lean{ThreeMajority.saturation_fail_le}
  \leanok
  \uses{lem:satcontract, def:satafter}
  Starting from a dissent bound $M$ ($500\log n\le M\le n/4$), offset $i$
  rounds in, the probability the dissent count exceeds
  $\mathrm{satAfter}(n,M,i+j)$ after $j$ further rounds is at most
  $j\,n^{-2}$, by induction on $j$ via the same $\mathrm{expList}$-
  conditioning idiom as Lemma~\ref{lem:growthfail}, with the tail direction
  flipped (dissent shrinking rather than opinion-$1$ count growing).
\end{lemma}

\begin{lemma}[Stage 2a, wrapped]\label{lem:stage2a}
  \lean{ThreeMajority.T2a, ThreeMajority.satAfter_quarter_le_floor, ThreeMajority.saturation_stage2a}
  \leanok
  \uses{lem:satfail}
  With $T_{2a}(n) := \ceil{6\log n}$: starting from any dissent bound
  $\le n/4$, after $T_{2a}(n)$ rounds the dissent count exceeds the floor
  $500\log n$ with probability at most $T_{2a}(n)\cdot n^{-2}$.
\end{lemma}

\begin{lemma}[Stage 2b]\label{lem:stage2b}
  \lean{ThreeMajority.saturation_stage2b}
  \leanok
  \uses{lem:satquad, lem:satclosed, lem:exppow, lem:loglog}
  Provided $\log n\ge30$: whenever $U_t\le500\log n$,
  $\Pr[U_{t+1}>10] \le \exp(-\log n) = 1/n$. One round, jumping from the
  $\Theta(\log n)$ floor directly to the fixed constant $10$: since the
  mean is already tiny ($\Theta((\log n)^2/n)$), this single round has an
  astronomically small failure probability once $n$ is large enough to beat
  the relevant polynomial in $\log n$, which a degree-$11$ instance of
  Lemma~\ref{lem:exppow}, combined with the tighter tangent-line bound of
  Lemma~\ref{lem:loglog}, already secures at this modest threshold.
\end{lemma}

\begin{lemma}[Stage 2c]\label{lem:stage2c}
  \lean{ThreeMajority.saturation_stage2c}
  \leanok
  \uses{lem:satquad}
  Whenever $U_t\le10$: $\Pr[U_{t+1}\ge1] \le 300/n$, via ordinary Markov
  applied to the nonnegative-integer $U_{t+1}$ ($\E[U_{t+1}]\le300/n$ by
  Lemma~\ref{lem:satquad} at $M=10$) --- no concentration needed for the
  very last step.
\end{lemma}

\section{Proof of the main theorem}\label{sec:proof-main}

\begin{lemma}[Stages 2b+2c combined]\label{lem:sat2b2c}
  \lean{ThreeMajority.ne_univ_iff_pos_dissent, ThreeMajority.saturation_2b2c}
  \leanok
  \uses{lem:stage2b, lem:stage2c}
  Two rounds, from $U_t\le500\log n$ to \emph{exactly} $U=0$, with total
  failure probability $\le\exp(-\log n)+300/n$, by chaining
  Lemma~\ref{lem:stage2b} then Lemma~\ref{lem:stage2c} via
  $\mathrm{expList}$-conditioning.
\end{lemma}

\begin{lemma}[All three saturation stages combined]\label{lem:sat2a2b2c}
  \lean{ThreeMajority.saturation_2a2b2c}
  \leanok
  \uses{lem:stage2a, lem:sat2b2c}
  $T_{2a}(n)+2$ rounds, from $U_t\le n/4$ to exactly $U=0$, with total
  failure probability $\le T_{2a}(n)\cdot n^{-2} + \exp(-\log n) + 300/n$.
\end{lemma}

\begin{theorem}[Main failure bound, raw form]\label{lem:mainfail}
  \lean{ThreeMajority.majority3_consensus_fail_le}
  \leanok
  \uses{lem:phase1, lem:sat2a2b2c}
  For $\log n\ge30$ and $|I_0|\ge\tfrac35n$: over $10+(T_{2a}(n)+2)$
  rounds,
  \[
    \Pr[I_T\ne V] \;\le\; 10e^{-n/10^7}
      + \bigl(T_{2a}(n)\,n^{-2} + e^{-\log n} + 300/n\bigr),
  \]
  by gluing Lemma~\ref{lem:phase1} (growth) to
  Lemma~\ref{lem:sat2a2b2c} (saturation), exactly as the push protocol's
  \texttt{prNotAllInformed\_le} glues its own two phases.
\end{theorem}

\begin{theorem}[Clean numeric bound]\label{lem:mainfailclean}
  \lean{ThreeMajority.hfloor4_of_hbig, ThreeMajority.majority3_consensus_fail_le_clean}
  \leanok
  \uses{lem:mainfail, lem:exppow}
  Under the same hypotheses, the four error terms of
  Theorem~\ref{lem:mainfail} are each bounded by $O(1/n)$ using
  Lemma~\ref{lem:exppow} (to show $n$ dominates the relevant polynomials in
  $\log n$ once $\log n\ge30$), giving the single clean bound
  $\Pr[I_T\ne V] \le 500/n$.
\end{theorem}

With $T = 10+(T_{2a}(n)+2)$, Theorem~\ref{thm:main}
(\texttt{majority3\_consensus\_whp}) follows from
Theorem~\ref{lem:mainfailclean} by
$\Pr[I_T=V] = 1-\Pr[I_T\ne V] \ge 1-500/n$.

\section{\texorpdfstring{3-Majority from any configuration in $O(n^{3/4}\log^{7/8} n)$ rounds}{3-Majority from any configuration}}\label{sec:bcekmn-any-start}

This section formalizes the upper bound for 3-Majority with any number of colours of
Berenbrink, Clementi, Elsässer, Kling, Mallmann-Trenn and Natale, \emph{Ignore or comply? On
breaking symmetry in consensus} (PODC 2017) \cite{BCEKMN17}: Theorem~4 of arXiv:1702.04921 (v1),
with the comparison framework of its Section~2 and its Lemmas~2 to~4
(\href{https://github.com/formal-dynamics/leanamics/blob/main/3-majority/FORMALIZATION_DIFFERENCES.md#3-majority-from-any-configuration-anystart-roadmap-maj-6-b}{differences
from the paper}).

\textbf{The second phase.} The proof of Theorem~4 applies Theorem~3.1 of Becchetti, Clementi,
Natale, Pasquale and Trevisan \cite{BCNPT16}, which is not part of \cite{BCEKMN17}.
Definition~\ref{def:as-phase2} states that result, Theorem~\ref{thm:as-main} takes it as a
hypothesis, exactly as the paper's proof uses it, and Section~\ref{sec:bcnpt16-many-opinions}
(roadmap MAJ-12~(a)) proves it (Theorem~\ref{thm:mo-phase2}), so that
Theorem~\ref{thm:as-main-unconditional}, Theorem~4 of \cite{BCEKMN17}, has no hypothesis.

\subsection{Model and majorization}

Agents $V=\{0,\dots,n-1\}$ hold colours in a finite set $\Sigma$; a configuration is a map
$c\colon V\to\Sigma$, with $c_a$ agents of colour $a$ and $\kappa(c)=|\{a : c_a>0\}|$ remaining
colours. Consensus is $\kappa(c)\le1$.

\begin{definition}[3-Majority and Voter with any number of colours]\label{def:as-model}
\lean{ThreeMajority.majColour,ThreeMajority.stepCol,ThreeMajority.runCol,ThreeMajority.voterStep,ThreeMajority.voterRun,ThreeMajority.colourCount,ThreeMajority.numColours}\leanok
\uses{def:round}
In a round of 3-Majority every agent draws an ordered sample triple $(s_1,s_2,s_3)$ uniformly
with replacement (Definition~\ref{def:round}); if $s_2$ and $s_3$ have the same colour it adopts
that colour, otherwise the colour of $s_1$. In a round of Voter every agent adopts the colour of
one uniformly random agent (itself included).
\end{definition}

The paper breaks a three-way tie by a uniformly random sample; since the samples are i.i.d.,
taking the first one gives the same law (Lemma~\ref{lem:as-alpha}).

\begin{lemma}[Two colours]\label{lem:as-bool}
\lean{ThreeMajority.stepCol_bool}\leanok\uses{def:as-model,def:step}
With colours $\{0,1\}$, a round of 3-Majority is the step of Definition~\ref{def:step}.
\end{lemma}
\begin{proof}\leanok
Case analysis on the colours of the three samples.
\end{proof}

\begin{definition}[Majorization, Schur-convex observables]\label{def:as-majorization}
\lean{ThreeMajority.Majorizes,ThreeMajority.countVec,ThreeMajority.SchurConvex}\leanok
\uses{def:as-model}
For $x,y\in\R^\Sigma$, $x\succeq y$ if $\sum x=\sum y$ and every sum of $y$ over a set
$S\subseteq\Sigma$ is at most a sum of $x$ over a set of the same size (equivalently, for every $j$
the $j$ largest entries of $y$ sum to at most the $j$ largest entries of $x$). An observable $\varphi$
of configurations is Schur-convex if $(c_a)_a\succeq(c'_a)_a$ implies $\varphi(c')\le\varphi(c)$.
\end{definition}

\begin{lemma}[Majorization and colours]\label{lem:as-colours}
\lean{ThreeMajority.numColours_le_of_majorizes,ThreeMajority.schurConvex_numColours_le}\leanok
\uses{def:as-majorization}
If $(c_a)_a\succeq(c'_a)_a$ then $\kappa(c)\le\kappa(c')$. Hence the indicator of
$\kappa(c)\le\kappa$ is Schur-convex.
\end{lemma}
\begin{proof}\leanok
The colours of $c'$ carry all $n$ agents, so a set of $\kappa(c')$ colours carries all agents of $c$.
\end{proof}

\subsection{Comparison of anonymous consensus processes}

\begin{definition}[AC-processes and dominance, Definitions~1 and~2]\label{def:as-ac}
\lean{ThreeMajority.acKernel,ThreeMajority.Dominates}\leanok\uses{def:as-majorization}
An anonymous consensus process with process function $\alpha$ moves from $c$ by letting every
agent draw its new colour independently from the probability vector $\alpha(c)$. The process
$\alpha$ dominates $\alpha'$ if $c\succeq c'$ implies $\alpha(c)\succeq\alpha'(c')$.
\end{definition}

\begin{proposition}[Proposition~1, Rinott]\label{prop:as-rinott}
\lean{ThreeMajority.multinomial_schurConvex,ThreeMajority.expect_le_of_transfer,ThreeMajority.expect_le_of_partial_sums,ThreeMajority.expect_le_of_majorizes}\leanok
\uses{def:as-majorization}
If $p\succeq q$ are probability vectors on $\Sigma$ and $\varphi$ is Schur-convex, the
expectation of $\varphi$ when the $n$ agents draw their colours i.i.d.\ from $q$ is at most the
expectation when they draw from $p$.
\end{proposition}
\begin{proof}\leanok
The paper cites this result without proof. For one Robin Hood transfer between colours $a$ and
$b$, the expectation, as a function of the mass $\theta$ of $a$ (with $a$ and $b$ sharing a fixed
total $s$), is nondecreasing on $[s/2,s]$: its derivative is a sum over agents, and swapping $a$
and $b$ for all agents and then for one agent pairs its terms into products of two factors with
the same sign. A chain of transfers (the Hardy--Littlewood--P\'olya construction) leads from $p$
to $q$ when $q$ is sorted and the partial sums of $p$ dominate those of $q$; sorting and
relabelling the colours, under which Schur-convex observables are invariant, reduce the general
case to this one.
\end{proof}

\begin{theorem}[Theorem~2, distributional form]\label{thm:as-comparison}
\lean{ThreeMajority.ac_comparison,ThreeMajority.ac_numColours}\leanok
\uses{def:as-ac,prop:as-rinott,lem:as-colours}
If $\alpha$ dominates $\alpha'$ and $c\succeq c'$, then for every Schur-convex $\varphi$ and every
$T$, $\E_{c'}[\varphi(P'_T)]\le\E_c[\varphi(P_T)]$. In particular, from the same configuration,
the dominating process has at most $\kappa$ colours at time $T$ with at least the probability of
the dominated one.
\end{theorem}
\begin{proof}\leanok
Induction on $T$. With $u_T$, $v_T$ the expectations at time $T$ for the dominated and the
dominating process, $w(y)=\max\{u_T(z):z\preceq y\}$ is Schur-convex and $u_T\le w\le v_T$ on
comparable configurations, so Proposition~\ref{prop:as-rinott} applied to $w$ gives the step
from $T$ to $T+1$. The second statement is the case of the indicator of
Lemma~\ref{lem:as-colours}. The paper's coupling (Lemma~1, from Strassen's theorem) is not
needed.
\end{proof}

\subsection{3-Majority is at least as fast as Voter}

\begin{lemma}[Equations~(1) and~(2)]\label{lem:as-alpha}
\lean{ThreeMajority.alpha3M,ThreeMajority.alphaVoter,ThreeMajority.alphaVoter_weight,ThreeMajority.alpha3M_weight,ThreeMajority.apply_ofStep_stepCol,ThreeMajority.apply_ofStep_voterStep}\leanok
\uses{def:as-model,def:as-ac}
Both processes are AC-processes. With $x=c/n$, Voter has $\alpha^V_a(c)=x_a$ and 3-Majority has
$\alpha^{3M}_a(c)=x_a(1+x_a-\|x\|_2^2)$.
\end{lemma}
\begin{proof}\leanok
The agents' samples are independent. For 3-Majority,
$[\mathrm{maj}=a]=[x_2=a][x_3=a]+[x_1=a]-[x_2=x_3][x_1=a]$, and the number of pairs of agents of
equal colour is $\sum_b c_b^2$.
\end{proof}

\begin{lemma}[Lemma~2]\label{lem:as-vs-voter}
\lean{ThreeMajority.dominates_alpha3M_alphaVoter,ThreeMajority.voter_le_threeMaj}\leanok
\uses{lem:as-alpha,thm:as-comparison}
3-Majority dominates Voter. Hence, from the same configuration, after $T$ rounds 3-Majority has
at most $\kappa$ colours with at least the probability that Voter has.
\end{lemma}
\begin{proof}\leanok
For a set $S$ of the largest entries of $x$, $\|x\|_2^2\sum_S x\le\sum_S x^2$, since
$\sum_{i\in S}\sum_j x_ix_j(x_i-x_j)\ge0$; this gives $\alpha^{3M}(c)\succeq\alpha^V(c)$, and with
$c\succeq c'$ also $\alpha^{3M}(c)\succeq\alpha^V(c')$. Conclude with
Theorem~\ref{thm:as-comparison}.
\end{proof}

\subsection{The Voter bound}

\begin{lemma}[Lemma~4 and Equation~(6), duality]\label{lem:as-dual}
\lean{ThreeMajority.walkStep,ThreeMajority.numColours_voterRun_le,ThreeMajority.voter_dual}\leanok
\uses{def:as-model}
Let coalescing random walks start on all $n$ nodes of the complete graph with self-loops. After
$T$ rounds, Voter has more than $k$ colours with at most the probability that more than $k$
walks remain.
\end{lemma}
\begin{proof}\leanok
Pathwise, the colours of Voter after rounds $y_1,\dots,y_T$ are colours of the nodes occupied
by the walks moved by $y_T,\dots,y_1$; the rounds are i.i.d., so reversing their order does not
change the law. The paper proves an equality under a coupling and uses only this inequality.
\end{proof}

\begin{lemma}[Equation~(7) and the expected number of walks]\label{lem:as-walks}
\lean{ThreeMajority.walk_drift,ThreeMajority.walk_drift_paper,ThreeMajority.walk_expect_le}\leanok
\uses{lem:as-dual}
If $X_t$ walks remain, $\E[X_{t+1}\mid X_t=x]\le x-x(x-1)/(3n)$ for every $x$, hence
$\le x-x^2/(10n)$ for $x\ge2$ (Equation~(7)). From any set of walks, $\E X_t\le1+3n/t$ for
$t\ge1$.
\end{lemma}
\begin{proof}\leanok
$\E[X_{t+1}\mid X_t=x]=n(1-(1-1/n)^x)$, the expected number of occupied cells when $x$ balls go
into $n$ cells; bound $(1-1/n)^x$ by the Bonferroni inequality of degree~3. The map
$x\mapsto x-x(x-1)/(3n)$ is concave and nondecreasing on $[0,n+1]$, so its tangent-line bound and
induction on $t$ give $\E X_t\le1+3n/t$. This replaces the paper's variable drift argument
for $\E[T^k_C]\le20n/k$.
\end{proof}

\begin{lemma}[Lemma~3]\label{lem:as-voter-whp}
\lean{ThreeMajority.voter_reduce_whp}\leanok\uses{lem:as-walks}
For $n\ge2$, $k\ge1$ and $T\ge24(n/k)\log n$, from any configuration Voter has more than $k$
colours after $T$ rounds with probability at most $1/n$.
\end{lemma}
\begin{proof}\leanok
Blocks of $\lceil6n/k\rceil$ rounds: by Markov's inequality on $X-1$ (the walks never vanish),
each block leaves more than $k$ walks with probability at most $1/2$, from any start. Restart
over $\lfloor T/\lceil6n/k\rceil\rfloor\ge\log_2 n$ blocks (for $k<n$; for $k\ge n$ there is
nothing to prove) and use Lemma~\ref{lem:as-dual}.
\end{proof}

\subsection{The main theorem}

\begin{lemma}[Phase~1]\label{lem:as-phase1}
\lean{ThreeMajority.threeMaj_reduce_whp}\leanok\uses{lem:as-vs-voter,lem:as-voter-whp}
For $n\ge2$, $k\ge1$ and $T\ge24(n/k)\log n$, from any configuration 3-Majority has more than
$k$ colours after $T$ rounds with probability at most $1/n$.
\end{lemma}
\begin{proof}\leanok
Lemma~\ref{lem:as-vs-voter} and Lemma~\ref{lem:as-voter-whp}.
\end{proof}

\begin{definition}[Phase~2: the cited result]\label{def:as-phase2}
\lean{ThreeMajority.Bcnpt16Phase2}\leanok\uses{def:as-model}
For $\varepsilon>0$, $\mathrm{Phase2}(\varepsilon)$ is the proposition (Theorem~3.1 of
\cite{BCNPT16}, Theorem~8 of \cite{BCEKMN17}): there are $C,N$ such that for $n\ge N$, from any
configuration with colours in $\{0,\dots,n-1\}$ and $k\le n^{1/3-\varepsilon}$ colours,
3-Majority has not reached consensus after any
$T\ge C(k^2\sqrt{\log n}+k\log n)(k+\log n)$ rounds with probability at most $1/n$.
This proposition is a hypothesis of Theorem~\ref{thm:as-main}; Theorem~\ref{thm:mo-phase2}
proves it for every $\varepsilon>0$ (Section~\ref{sec:bcnpt16-many-opinions}). The cited paper's
``w.h.p.''\ is $1-O(n^{-\lambda})$ for some $\lambda>0$, which gives the form $1/n$ by repeating
the phase a constant number of times; the proof of Theorem~\ref{thm:mo-phase2} gives the form
$1/n$ directly.
\end{definition}

\begin{lemma}[The asymptotics of Theorem~4]\label{lem:as-asymptotics}
\lean{ThreeMajority.phase1Colours,ThreeMajority.anyStart_asymptotics}\leanok
For every $A\ge0$ there are $C,N$ such that for $n\ge N$ and
$k=\lfloor n^{1/4}\log^{1/8}n\rfloor$: $n\ge2$, $1\le k\le n^{1/3-1/24}$, and
$24(n/k)\log n+1+A(k^2\sqrt{\log n}+k\log n)(k+\log n)\le Cn^{3/4}\log^{7/8}n$.
\end{lemma}
\begin{proof}\leanok
From $\log y\le y^{1/4}$ and $\log^{1/8}y\le y^{1/24}$ for $y$ large.
\end{proof}

\begin{theorem}[Theorem~4 from Phase~2]\label{thm:as-main}
\lean{ThreeMajority.threeMaj_anyStart_consensus}\leanok
\uses{lem:as-phase1,def:as-phase2,lem:as-asymptotics}
Assume $\mathrm{Phase2}(1/24)$. There are $C,N$ such that for $n\ge N$, from any configuration
with colours in $\{0,\dots,n-1\}$, 3-Majority has not reached consensus after any
$T\ge Cn^{3/4}\log^{7/8}n$ rounds with probability at most $2/n$.
\end{theorem}
\begin{proof}\leanok
Let $k=\lfloor n^{1/4}\log^{1/8}n\rfloor$ and $T_1=\lceil24(n/k)\log n\rceil$. After $T_1$ rounds,
more than $k$ colours remain with probability at most $1/n$ (Lemma~\ref{lem:as-phase1});
otherwise the current configuration has at most $k\le n^{1/3-1/24}$ colours, and
$\mathrm{Phase2}(1/24)$ bounds the failure of the remaining $T-T_1$ rounds by $1/n$.
Lemma~\ref{lem:as-asymptotics} fits both phases in $Cn^{3/4}\log^{7/8}n$ rounds. The proof in the
paper needs a minor correction: it applies Lemma~3 with $k=n^{1/4}$, which gives
$O(n^{3/4}\log n)$ rounds for Phase~1; the choice $k=n^{1/4}\log^{1/8}n$ balances the two phases
and gives the bound as stated.
\end{proof}

\begin{theorem}[Theorem~4]\label{thm:as-main-unconditional}
\lean{ThreeMajority.threeMaj_anyStart_consensus_unconditional}\leanok
\uses{thm:as-main,thm:mo-phase2}
There are $C,N$ such that for $n\ge N$, from any configuration with colours in
$\{0,\dots,n-1\}$, 3-Majority has not reached consensus after any $T\ge Cn^{3/4}\log^{7/8}n$
rounds with probability at most $2/n$.
\end{theorem}
\begin{proof}\leanok
Theorem~\ref{thm:as-main}, whose hypothesis $\mathrm{Phase2}(1/24)$ is
Theorem~\ref{thm:mo-phase2} for $\varepsilon=1/24$.
\end{proof}

\section{3-Majority with many opinions}\label{sec:bcnpt16-many-opinions}

This section formalizes the adversary-free case of Becchetti, Clementi, Natale, Pasquale and
Trevisan, \emph{Stabilizing consensus with many opinions} (SODA 2016) \cite{BCNPT16}: its
Theorem~3.1, with its Lemmas~2.1 and~3.2 to~3.5 and Lemma~B.1, in the numbering of
arXiv:1508.06782 (v1). The
\href{https://github.com/formal-dynamics/leanamics/blob/main/3-majority/FORMALIZATION_DIFFERENCES.md#3-majority-with-many-opinions-manyopinions-roadmap-maj-12-a}{differences
from the paper} are listed in the repository. Theorem~3.1 is the proposition $\mathrm{Phase2}(\varepsilon)$ of
Definition~\ref{def:as-phase2} (Theorem~\ref{thm:mo-phase2}), so it makes Theorem~4 of
\cite{BCEKMN17} unconditional (Theorem~\ref{thm:as-main-unconditional}).

The process is the 3-Majority of Definition~\ref{def:as-model}, whose rule is, pointwise, the rule
of the paper (three samples uniform with repetition; on three different colours, the first one),
with colours in $\{0,\dots,n-1\}$. We write $c_a$ for the support of colour $a$, $C'_a$ for its
support after one round, $j$ or $k$ for the number of colours, and $m=16\sqrt{jn\log n}$. Where
the paper assumes $j\le n^{1/3-\varepsilon}$, the statements below assume $Cj^3\log n\le n$ for a
constant $C$, which is what the proofs use; it holds for $j\le n^{1/3-\varepsilon}$ and $n$
large. The proofs follow the paper's, with the deviations and the minor corrections noted in each
proof.

\subsection{One round}

\begin{lemma}[Lemma~2.1]\label{lem:mo-expect}
\lean{ThreeMajority.numColours_eq_card_image,ThreeMajority.colourCount_eq_card_filter,ThreeMajority.stepCol_apply,ThreeMajority.adoptInd,ThreeMajority.colourCount_stepCol_eq_sum,ThreeMajority.avg_colourCount_stepCol,ThreeMajority.avg_colourCount_stepCol_le}\leanok
\uses{def:as-model}
For any finite set of colours, any configuration $c$ and any colour $a$,
\[
  \E[C'_a]=c_a\Bigl(1+\frac{c_a}{n}-\sum_b\frac{c_b^2}{n^2}\Bigr)
  \le c_a\Bigl(1+\frac{c_a}{n}-\frac{1}{\kappa(c)}\Bigr).
\]
\end{lemma}
\begin{proof}\leanok
$C'_a$ is the sum over the agents of independent indicators of adopting $a$, each with mean
$x_a(1+x_a-\|x\|_2^2)$ for $x=c/n$ (as in Lemma~\ref{lem:as-alpha}). The inequality is the
Cauchy--Schwarz bound $\sum_b c_b^2\ge n^2/\kappa(c)$.
\end{proof}

\begin{lemma}[One-round tails and anti-concentration]\label{lem:mo-tails}
\lean{ThreeMajority.avg_colourCount_stepCol_ge,ThreeMajority.avg_colourCount_stepCol_le_tail,Dynamics.avg_pz_zero_one,ThreeMajority.avg_abs_colourCount_stepCol_sub_ge,ThreeMajority.sum_variance_adoptInd_ge}\leanok
\uses{lem:mo-expect,lem:chernoff,lem:indep}
Let $\sigma_a^2$ be the sum of the variances of the agents' indicators of adopting $a$.
\begin{enumerate}
  \item If $\E[C'_a]\le\mu<k$, then $\Pr(C'_a\ge k)\le\exp(-(k-\mu)^2/(k+\mu))$; if
    $0<k<\mu\le\E[C'_a]$, then $\Pr(C'_a\le k)\le\exp(-(\mu-k)^2/(2\mu))$.
  \item If $\sigma_a^2\ge1$, then $\E|C'_a-\theta|\ge(9/256)\,\sigma_a$ for every $\theta\in\R$.
  \item If $c$ has $j\ge2$ colours, all with more than $n/j-nu$ agents, where $u\ge0$ and
    $ju\le1/250$, then $\sigma_a^2\ge n/(4j)$ for every colour $a$ of $c$.
\end{enumerate}
\end{lemma}
\begin{proof}\leanok
(1) The Chernoff bounds of Section~\ref{sec:chernoff} for a sum of independent indicators.
(2) If $|\E[C'_a]-\theta|\ge\sigma_a/4$, this is Jensen's inequality. Otherwise, by the
Paley--Zygmund inequality (from the second and fourth moments of a sum of independent centred
$\{0,1\}$-coordinates) the centred sum $S$ satisfies $\sigma_a^2\le4S^2$ with probability at least
$9/64$, and there $|C'_a-\theta|\ge\sigma_a/2-\sigma_a/4$. (3) The adoption probability $p$ of
Lemma~\ref{lem:mo-expect} lies between $(1-1/250)^2/j$ and $(1/2+1/250)(1+1/250)$, and
$\sigma_a^2=np(1-p)$.
\end{proof}

\subsection{Hitting times under a drift}

\begin{lemma}[Lemma~3.2, finite horizon]\label{lem:mo-drift}
\lean{ThreeMajority.hitProb_additive_drift}\leanok
Let $K$ be a Markov kernel on a finite set, $I$ a set with $K(x,I)=1$ for $x\in I$, $B$ a target
set with hitting time $\tau$, and $f$ a real function. Let $\beta,\eta\ge0$, and assume that
$f(x)\le M$ for $x\in I$ and that, for $x\in I\setminus B$, $f(x)\le\beta$,
$\E[f(X_1)\mid X_0=x]\ge f(x)+\lambda$ and $\Pr(f(X_1)>\beta\mid X_0=x)\le\eta$. Then for every
$a\in I$ with $f(a)\ge0$ and every $t$, $\lambda t\Pr_a(\tau>t)\le\beta+M\eta t$.
\end{lemma}
\begin{proof}\leanok
Let $X'$ be the chain stopped at $B$. Summing the drift over the steps before $\tau$,
$\E_a[f(X'_t)]\ge f(a)+\lambda\E_a[\min(t,\tau)]\ge\lambda t\Pr_a(\tau>t)$. By a union bound over
the steps, $f(X'_t)>\beta$ with probability at most $\eta t$, and $f\le M$, so
$\E_a[f(X'_t)]\le\beta+M\eta t$. The paper bounds $\E[\tau]$ by optional stopping, under a
hypothesis on the overshoot at time $\tau$ (Lemma~B.1); the finite horizon and the one-step
overshoot bound avoid both.
\end{proof}

\begin{lemma}[Hit, then absorb]\label{lem:mo-strong-markov}
\lean{ThreeMajority.hitProb_mul_le_event}\leanok
Let $P$ be an absorbing set for $K$. If from every state of $B$ the chain is in $P$ after $s$
steps with probability at least $q$, then from every state $a$ it is in $P$ at time $t+s$ with
probability at least $q\Pr_a(\tau\le t)$.
\end{lemma}
\begin{proof}\leanok
The strong Markov property at $\tau$, proved by induction on $t$ without stopping times: from a
state of $B$, the hypothesis and absorption; otherwise, condition on the first step.
\end{proof}

\subsection{Symmetry breaking}

\begin{definition}[Broken symmetry]\label{def:mo-broken}
\lean{ThreeMajority.BrokenSymmetry,ThreeMajority.brokenSymmetry_iff}\leanok
\uses{def:as-model}
A configuration $x$ has broken symmetry for $j$ colours if $\kappa(x)<j$, or $\kappa(x)\le j$ and
some colour has between $1$ and $n/j-16\sqrt{jn\log n}$ agents.
\end{definition}

The paper's threshold is $n/j-\sqrt{jn\log n}$; the factor $16$ is needed in
Lemma~\ref{lem:mo-stage1}. The death of a colour also breaks the symmetry (the paper's minimum
is over the initial colours, so a dead colour has support $0$).

\begin{lemma}[Equation~(3)]\label{lem:mo-min-drift}
\lean{ThreeMajority.min_count_drift}\leanok
\uses{lem:mo-expect,lem:mo-tails}
There are $\delta>0$ and $C$ such that, for every configuration $x$ with $j=\kappa(x)\ge2$
colours, $Cj^3\log n\le n$ and every colour of $x$ above $n/j-m$ agents, the smallest support
among the colours of $x$ (a colour that dies counting $0$) decreases in expectation by at least
$\delta\sqrt n/(j\sqrt j)$ in one round.
\end{lemma}
\begin{proof}\leanok
Take $\delta=1/500$, $C=1024$ and $w=\sqrt{n/j}$. If the smallest support $c_{\min}$ is at most
$n/j-w/250$, Lemma~\ref{lem:mo-expect} for a smallest colour, with $c_{\min}\ge n/(2j)$
(Equation~(5)), gives the drift. Otherwise all colours are above $n/j-w/250$. The supports sum to
$n$, so the smallest new support is at most $n/j-\frac1{2j}\sum_a|C'_a-n/j|$, and
$\E|C'_a-n/j|\ge(9/256)\sigma_a\ge(9/512)w$ by Lemma~\ref{lem:mo-tails}; hence the expected
smallest new support is at most $n/j-9w/1024<c_{\min}-w/1000$. The paper argues in this case with
the Central Limit Theorem and negative association instead.
\end{proof}

\begin{lemma}[Equation~(6) and Lemma~B.1, one step]\label{lem:mo-overshoot}
\lean{ThreeMajority.min_count_overshoot}\leanok
\uses{lem:mo-expect,lem:mo-tails}
There is $C$ such that, under the hypotheses of Lemma~\ref{lem:mo-min-drift}, one round puts some
colour of $x$ below $n/j-64\sqrt{jn\log n}$ agents with probability at most $1/n^2$.
\end{lemma}
\begin{proof}\leanok
Take $C=2^{14}$. Every support is at most $n/j+(j-1)m$, so by Lemma~\ref{lem:mo-expect} every
colour of $x$ has $\E[C'_a]\ge n/j-2m$. The Chernoff lower tail at $n/j-4m$
(Lemma~\ref{lem:mo-tails}) and a union bound over the at most $n$ colours give $1/n^2$. The paper
bounds the overshoot at the hitting time instead (Lemma~B.1); its proof needs a minor correction,
since it treats the smallest support as a Markov chain, while it is a function of the
configuration. The one-step form avoids this argument.
\end{proof}

\begin{lemma}[Lemma~3.3, symmetry breaking]\label{lem:mo-breaking}
\lean{ThreeMajority.symmetry_breaking,ThreeMajority.breakPotential}\leanok
\uses{def:mo-broken,lem:mo-drift,lem:mo-min-drift,lem:mo-overshoot}
There is $C$ such that, from every configuration $c$ with $j\ge2$ colours and $Cj^3\log n\le n$,
3-Majority reaches broken symmetry for $j$ colours within any $t\ge Cj^2\sqrt{\log n}$ rounds with
probability at least $1/2$.
\end{lemma}
\begin{proof}\leanok
Lemma~\ref{lem:mo-drift} with $I=\{\kappa\le j\}$, $B$ the configurations with broken symmetry,
$f=n/j$ minus the smallest support among the $j$ colours (the paper's $Y_t$),
$\lambda=\delta\sqrt n/(j\sqrt j)$ (Lemma~\ref{lem:mo-min-drift}), $\beta=64\sqrt{jn\log n}$ and
$\eta=1/n^2$ (Lemma~\ref{lem:mo-overshoot}), and $M=n/j$. For $t\ge Cj^2\sqrt{\log n}$ (the proof
takes $C=145\,409$), $\lambda t\ge4\beta$ and $M\eta t\le\lambda t/4$, so $\Pr(\tau>t)\le1/2$.

The paper states the conclusion at the fixed time $t$, while its proof bounds the hitting time,
which is what is stated here; its composition in Theorem~3.1 needs a minor correction, the strong
Markov property at the hitting time (Lemma~\ref{lem:mo-strong-markov}). The statement also needs
a minor correction: it needs the hypothesis $j^3\log n=O(n)$, as do the paper's Lemmas~3.4
and~3.5 (otherwise
the threshold $n/j-\sqrt{jn\log n}$ is negative and the event has probability $0$). Finally, its
bound $\E[\tau]<j^2\sqrt{\log n}$ should read $\E[\tau]\le(2\alpha/\varepsilon)j^2\sqrt{\log n}$.
\end{proof}

\subsection{Dropping}

\begin{lemma}[Lemma~3.4, dropping stage~1]\label{lem:mo-stage1}
\lean{ThreeMajority.dropping_stage1,ThreeMajority.stage1Level,ThreeMajority.stage1_step}\leanok
\uses{lem:mo-expect,lem:mo-tails}
There is $C$ such that, for $j\ge2$ with $Cj^3\log n\le n$, from every configuration with at most
$j$ colours and every colour $a$ with at most $n/j-16\sqrt{jn\log n}$ agents, at every time $t$
with $Cj\log n\le t\le n$ colour $a$ has more than $Cj^2\log n$ agents with probability at most
$1/n$.
\end{lemma}
\begin{proof}\leanok
Take $C=240$. As in Appendix~B of the paper, the gap $n/j-c_a$ grows by a factor $1+1/(4j)$ per
round until $c_a\le n/(2j)$, and then $c_a$ shrinks by a factor $1-1/(4j)$ per round down to a
floor. Here both stages are one moving level,
$\max\{n/j-m(1+1/(4j))^s,\ (n/(2j))(1-1/(4j))^{s-s_A},\ K\}$ at round $s$, with
$K=\min\{240j^2\log n,n/(2j)\}$ and $s_A=\lceil5j\log n\rceil$. By Lemma~\ref{lem:mo-expect} and
the Chernoff upper tail (Lemma~\ref{lem:mo-tails}), $c_a$ exceeds the level of the next round
with probability at most $1/n^2$; after $s_A+\lceil4j\log n\rceil\le240j\log n$ rounds the level
is $K$, and a union bound over the $t\le n$ rounds gives $1/n$.

The proof in the paper needs a minor correction. At its threshold $n/j-\sqrt{jn\log n}$, and
again near its floor $j^2\log n$, the one-round bound of its Equation~(7) is only about
$n^{-1/12}$, and the union bound over the rounds exceeds $1$ for large $j$ (for instance
$j\ge20n^{1/10}$, which $j\le n^{1/3-\varepsilon}$ allows). The threshold $16\sqrt{jn\log n}$,
which Lemma~\ref{lem:mo-breaking} reaches within the same $O(j^2\sqrt{\log n})$ rounds, and the
floor $240j^2\log n$ make each round fail with probability at most $1/n^2$.
\end{proof}

\begin{lemma}[Lemma~3.5, dropping stage~2]\label{lem:mo-stage2}
\lean{ThreeMajority.dropping_stage2,ThreeMajority.avg_min_colourCount_stepCol_le,ThreeMajority.expList_min_colourCount_le}\leanok
\uses{lem:mo-expect,lem:mo-tails}
There is $C$ such that, for $j\ge2$ with $Cj^3\log n\le n$, from every configuration with at most
$j$ colours and every colour $a$ with at most $n/(2j)$ agents, colour $a$ has died after any
$t\ge Cj\log n$ rounds with probability at least $1/2$.
\end{lemma}
\begin{proof}\leanok
Take $C=1024$. As in the paper's sketch, with the potential $\psi=\min\{c_a,n/(2j)\}$:
Lemma~\ref{lem:mo-expect} gives $\E[\psi']\le(1-1/(2j))\psi+(n/(4j^2))[c_a>n/(2j)]$, and by the
Chernoff upper tail $c_a$ escapes above $n/(2j)$ with probability at most $1/n^4$ per round.
Hence $\E[\psi_s]\le(1-1/(2j))^s\,n/(2j)+(n/(4j^2))\,n^{-4}s^2$, which is at most $1/2$ at
$s=\lceil2j\log n\rceil$, and Markov's inequality gives $\Pr(c_a\ge1)\le1/2$. Later times follow
since a dead colour stays dead.
\end{proof}

\subsection{Consensus}

\begin{lemma}[One phase]\label{lem:mo-phase}
\lean{ThreeMajority.phase_colour_dies}\leanok
\uses{lem:mo-breaking,lem:mo-stage1,lem:mo-stage2,lem:mo-strong-markov}
There is $C$ such that, from every configuration with $j\ge2$ colours and $Cj^3\log n\le n$,
after any $t\ge C(j^2\sqrt{\log n}+j\log n)$ rounds fewer than $j$ colours remain with probability
at least $1/5$.
\end{lemma}
\begin{proof}\leanok
From a configuration with broken symmetry, either fewer than $j$ colours remain, or a colour has
at most $n/j-16\sqrt{jn\log n}$ agents; then Lemma~\ref{lem:mo-stage1} brings it to at most
$240j^2\log n\le n/(2j)$ agents except with probability $1/n$, and Lemma~\ref{lem:mo-stage2} kills
it with probability at least $1/2$: in total at least $\frac12(1-\frac1n)\ge\frac25$. Fewer than
$j$ colours is absorbing, so Lemma~\ref{lem:mo-strong-markov} at the hitting time of broken
symmetry and Lemma~\ref{lem:mo-breaking} give $\frac12\cdot\frac25=\frac15$ (the proof takes
$C=293\,349$). The paper's proof of Theorem~3.1 needs a minor correction here: its probability
$1/4$ ignores the failure probability of Lemma~3.4; any positive constant suffices.
\end{proof}

\begin{lemma}[Amplification over phases]\label{lem:mo-phases}
\lean{ThreeMajority.consensus_of_phases}\leanok
\uses{def:as-model}
If from every configuration with between $2$ and $k$ colours a phase of $L$ rounds removes a
colour with probability at least $p\le1$, then from every configuration with at most $k$ colours
consensus fails after $M$ phases with probability at most $(1-p/2)^M2^k$.
\end{lemma}
\begin{proof}\leanok
The potential equal to $2^{\kappa}$ without consensus and to $0$ at consensus: colours never
increase, and with probability at least $p$ one of them disappears, so over one phase its
expectation is multiplied by at most $(1-p)+p/2=1-p/2$. This makes precise the paper's
``standard probabilistic arguments''.
\end{proof}

\begin{theorem}[Theorem~3.1, for $Ck^3\log n\le n$]\label{thm:mo-main}
\lean{ThreeMajority.consensus_many_opinions,ThreeMajority.expList_not_consensus_le}\leanok
\uses{lem:mo-phase,lem:mo-phases}
There is $C$ such that, from every configuration with $k$ colours and $Ck^3\log n\le n$,
3-Majority has not reached consensus after any $T\ge C(k^2\sqrt{\log n}+k\log n)(k+\log n)$
rounds with probability at most $1/n$.
\end{theorem}
\begin{proof}\leanok
Phases of $L=\lceil C_1(k^2\sqrt{\log n}+k\log n)\rceil$ rounds, with $C_1$ the constant of
Lemma~\ref{lem:mo-phase}, which applies to every configuration with $2\le j\le k$ colours. By
Lemma~\ref{lem:mo-phases} with $p=1/5$ and $M=10(k+\lceil\log n\rceil)$ phases, consensus fails
after $ML$ rounds with probability at most $(9/10)^M2^k\le1/n$, and consensus is absorbing. The
proof takes $C=15(C_1+1)=4\,400\,250$.
\end{proof}

\begin{theorem}[Theorem~3.1]\label{thm:mo-phase2}
\lean{ThreeMajority.bcnpt16Phase2_iff,ThreeMajority.bcnpt16Phase2_of_pos}\leanok
\uses{thm:mo-main,def:as-phase2}
For every $\varepsilon>0$, $\mathrm{Phase2}(\varepsilon)$ holds: there are $C,N$ such that for
$n\ge N$, from every configuration with colours in $\{0,\dots,n-1\}$ and $k\le n^{1/3-\varepsilon}$
colours, 3-Majority has not reached consensus after any
$T\ge C(k^2\sqrt{\log n}+k\log n)(k+\log n)$ rounds with probability at most $1/n$.
\end{theorem}
\begin{proof}\leanok
Theorem~\ref{thm:mo-main} with the same $C$: $k^3\le n^{1-3\varepsilon}$ and
$|C|\log n\le n^{3\varepsilon}$ for $n\ge N$, so $Ck^3\log n\le n$.
\end{proof}

With $\varepsilon=1/24$, Theorem~\ref{thm:mo-phase2} is the hypothesis of
Theorem~\ref{thm:as-main}, which gives Theorem~\ref{thm:as-main-unconditional}.

\section{The Lean formalization}\label{sec:lean}

The argument above is fully formalized (no \texttt{sorry}) in the
accompanying repository (library \texttt{ThreeMajority}, namespace
\texttt{ThreeMajority}). File structure (dependency order):
\texttt{Prob.lean} (Section~\ref{sec:prelim}: aliases of the finite-uniform
\texttt{avg}/\texttt{expList} machinery of the shared \texttt{dynamics}
library, including the independence toolkit \texttt{avg\_prod\_pi} and
friends), \texttt{Model.lean}
(Section~\ref{sec:model}), \texttt{Bounds.lean} (Section~\ref{sec:ineq}),
the shared \texttt{Dynamics/Chernoff.lean} (Section~\ref{sec:chernoff}),
\texttt{OneRound.lean}
(Section~\ref{sec:map}), \texttt{Growth.lean} (Section~\ref{sec:growth}),
\texttt{Saturation.lean} (Section~\ref{sec:saturation}), \texttt{Main.lean}
(Section~\ref{sec:proof-main} and the main theorem). The files
\texttt{AnyStartModel}, \texttt{AnyStartMajorization}, \texttt{AnyStartRinott},
\texttt{AnyStartComparison}, \texttt{AnyStartVsVoter}, \texttt{AnyStartVoter},
\texttt{AnyStartAsymptotics} and \texttt{AnyStartMain} formalize
Section~\ref{sec:bcekmn-any-start}, on the shared finite \texttt{Distribution} and
\texttt{Kernel} layer of \texttt{dynamics}; the files \texttt{ManyOpinionsExpect},
\texttt{ManyOpinionsDrift}, \texttt{ManyOpinionsTails},
\texttt{ManyOpinionsBreaking}, \texttt{ManyOpinionsDropping} and \texttt{ManyOpinionsMain}
formalize Section~\ref{sec:bcnpt16-many-opinions}, with the moment bounds of
\texttt{Moments} in \texttt{dynamics}. Design constraint
preserved from \texttt{rumor\_spread}: no measure theory, no
\texttt{PMF}/\texttt{ENNReal}, no Mathlib \texttt{ProbabilityTheory}
kernels --- everything is finite sums, plus \texttt{Real.exp}/\texttt{Real.log}
where the Chernoff bound genuinely needs them (structurally central here,
unlike the push protocol, where \texttt{Real.exp} appears only in a numeric
side lemma).

\begin{thebibliography}{9}
\bibitem{DGMSS11} B.~Doerr, L.A.~Goldberg, L.~Minder, T.~Sauerwald,
C.~Scheideler. \emph{Stabilizing consensus with the power of two choices}.
SPAA 2011.
\bibitem{BCNPS15} L.~Becchetti, A.~Clementi, E.~Natale, F.~Pasquale,
R.~Silvestri. \emph{Plurality consensus in the gossip model}. SODA 2015.
\bibitem{BCEKMN17} P.~Berenbrink, A.~Clementi, R.~Elsässer, P.~Kling,
F.~Mallmann-Trenn, E.~Natale. \emph{Ignore or comply? On breaking symmetry in consensus}.
PODC 2017. arXiv:1702.04921.
\bibitem{BCNPT16} L.~Becchetti, A.~Clementi, E.~Natale, F.~Pasquale, L.~Trevisan.
\emph{Stabilizing consensus with many opinions}. SODA 2016. arXiv:1508.06782.
\end{thebibliography}
